%=======================================================================
%  Radioprotection (EDP Sciences / SFRP) submission
%  Re-aimed from the rejected Space Weather manuscript 1072911.
%  v3 (major revision): ALL dose values are real CARI-7 v4.1.3 executions
%   - Altitude base: log-quadratic in ln D (replaces single exponential);
%     MARE 0.10% over FL200-FL400, 1.05% over FL100-FL400
%   - Latitude factor: direct CARI-7 runs at 11 geographic nodes
%     (VCR 0.0-17.2 GV) + 3 held-out nodes (out-of-sample MARE 1.9%)
%   - Solar scaling and annual doses from real CARI-7 (HCP 1141 / 275 MV)
%   - Machine-learning tier removed; physics corrected (SC25 max Oct 2024;
%     GLE 77 acknowledged; H*(10)/E over-claim removed)
%=======================================================================
\documentclass[11pt,a4paper]{article}
\usepackage[utf8]{inputenc}
\usepackage[T1]{fontenc}
\usepackage[margin=2.5cm]{geometry}
\usepackage{amsmath,amssymb}
\usepackage{graphicx}
\usepackage{epstopdf}
\usepackage{textcomp}
\usepackage{tabularx}
\usepackage[round,authoryear]{natbib}
\usepackage{lineno}
\usepackage{setspace}
\usepackage{hyperref}
\onehalfspacing
\linenumbers
\setlength{\emergencystretch}{3em}
\newcommand{\uSvh}{\,\textmu Sv\,h$^{-1}$}
\newcommand{\mSvy}{\,mSv\,yr$^{-1}$}

\title{\bfseries A minimum-data altitude--latitude formula for rapid
screening of aircrew cosmic-radiation dose, derived from CARI-7}

\author{%
Doo-Hee Lee\textsuperscript{1}, Jung-Mi Park\textsuperscript{1},
Ho Sub Chi\textsuperscript{1}, Jin Yoo\textsuperscript{1},
Myeongsik Shin\textsuperscript{1}, Hyunsung Kim\textsuperscript{1},\\
Jun Yoo\textsuperscript{1}, Dongyoul Lee\textsuperscript{2},
Ku Kang\textsuperscript{1,$*$}}
\date{}

\begin{document}
\maketitle

\noindent\textsuperscript{1}CBRN Defense Research Institute, Ministry of
National Defense, Seoul 06796, Republic of Korea.\\
\textsuperscript{2}Department of Physics and Chemistry, Korea Military
Academy, Seoul 01805, Republic of Korea.\\
\textsuperscript{$*$}Corresponding author: bisu9082@gmail.com

%-----------------------------------------------------------------------
\begin{abstract}
\noindent
Aircrew are occupationally exposed, yet screening which profiles need
monitoring still relies on detailed route-based codes. We derive a
minimum-data screening formula for cosmic-radiation effective dose rate
from the FAA CARI-7 model. Over the FL200--FL400 cruise band the dose rate
follows a compact log-quadratic in altitude,
$\ln D=-0.0232\,h_{\rm km}^{2}+0.701\,h_{\rm km}-4.172$, with a mean
absolute relative error (MARE) of 0.10\% against CARI-7 (1.05\% over
FL100--FL400), from altitude alone. A saturation-type geomagnetic
correction, $f_{\rm lat}(R_c)=(1+aR_{c,\rm ref}^{\,n})/(1+aR_c^{\,n})$
($a=0.024$, $n=1.52$), was fitted to CARI-7 rates at eleven nodes spanning
$R_c=0$--17\,GV and reproduces three held-out nodes with an out-of-sample
MARE of 1.9\%. All values are real CARI-7 (v4.1.3) executions; the formula
reproduces CARI-7 and is not an independent measurement validation. Annual
effective doses near 40{,}000\,ft reach 1.7--2.2\,mSv\,yr$^{-1}$, and routes
with $R_c\lesssim7.3$\,GV at FL300 cross the 1\,mSv\,yr$^{-1}$ public
reference level. On five real routes the screen reproduces CARI-7
full-flight doses to within 8\% (long-haul) and matches its
1\,mSv\,yr$^{-1}$ flag. It is a pocket-calculator complement to full
route-based dosimetry.
\end{abstract}

\vspace{0.4em}
\noindent\textbf{Keywords:} \textbf{cosmic radiation} / \textbf{aircrew
dosimetry} / \textbf{CARI-7} / \textbf{geomagnetic cutoff rigidity} /
\textbf{dose screening}

%-----------------------------------------------------------------------
\section{Introduction}

Galactic cosmic radiation (GCR) is a pervasive radiological hazard at
aviation altitudes, where the shielding column of overlying atmosphere is
substantially reduced \citep{obrien1997atmospheric,icrp2007}; the same
primary flux drives atmospheric ionization and cosmogenic-nuclide
production \citep{makrantoni2022,environ}. At cruising altitudes the
effective dose rate exceeds sea-level background by one to three orders of
magnitude \citep{pelliccioni2000,ferrari1997,lindborg2004}, and
international bodies classify aircrew as an occupationally exposed
population requiring systematic dose assessment
\citep{icrp2007,icrp1991,icrp132}. Epidemiological cohorts of aircrew
sustain this concern \citep{hammer2009,grajewski2011,dreger2020}. In the
European Union, Council Directive 2013/59/Euratom (Article 35) requires
assessment of aircrew exposure where the effective dose is liable to exceed
1\,mSv\,yr$^{-1}$ \citep{eudirective2013}. The same secondary field also
drives single-event effects in avionics and microelectronics as device
sizes shrink \citep{normand1996,dyerqinetiq2007}.

Operational dose assessment is well served by physics-based route codes:
the FAA CARI-7 suite \citep{anic2016,copeland2008}, the European EPCARD and
AVIDOS systems \citep{latocha2016}, the French SIEVERT operational service
\citep{sievert,clairand2009sievert,bottollier2006}, the NAIRAS real-time
model \citep{mertens2013}, and the PHITS-based PARMA/EXPACS model
\citep{sato2015parma}. Analytical and measurement-based representations of
the global aviation dose field have also been developed from the NAIRAS
model and the ARMAS in-flight database
\citep{tobiska2018,meier2018,meier2014swindex,phoenix2026}. These tools and
representations are accurate but require route geometry, date, and, in most
cases, software or network access.

A distinct and practically important need is \emph{rapid, minimum-data
screening}: a first-pass estimate of which flight profiles may warrant
closer monitoring, obtainable from flight altitude alone, for crew
briefing, retrospective analysis from flight logs, or resource-limited
settings. Compact dose terms are already used in our group's downstream
frameworks---offline network-denied dose assessment \citep{resil} and an
aircrew cochlear-risk model with a CARI-7-derived dose term
\citep{lee2026cochlear}---motivating a documented screening formula. The increase of dose rate with
altitude is well established \citep{lindborg2004,schraube1998}, and compact
altitude--dose relations have been reported for particular latitude bands
and national airspaces
\citep{vukovic2022,bartlett2004,zarrouk2008,makrantoni2022}. Our aim is not
to introduce a new dose
model but to provide a transparent, pocket-calculator screening
approximation that is tied explicitly to a recognised reference (CARI-7)
and that carries an honest statement of its error and validity range.
Existing compact forms meet only part of this need: operational services
such as SIEVERT are not reducible to a hand calculation, published
analytical representations of the dose field still require software
\citep{tobiska2018,meier2018}, and reported exponential altitude relations
are single-band and are not accompanied by an out-of-sample error budget
\citep{vukovic2022,bartlett2004}. To our knowledge, no prior work provides
a single hand-evaluable expression, anchored to a named reference code and
validated out-of-sample, that spans both altitude and geomagnetic
latitude.

The objectives of this study are: (i) to derive, from CARI-7, a compact
closed-form altitude expression for the FL200--FL400 cruise band and
quantify how closely it reproduces the reference; (ii) to add a
geomagnetic cutoff-rigidity correction, fitted and independently tested
against CARI-7, so the screen applies from polar to equatorial routes;
(iii) to use the resulting formula to flag flight profiles whose annual
dose approaches or exceeds the 1\,mSv\,yr$^{-1}$ public reference level; and
(iv) to characterise, from the same CARI-7 runs, the structure that such a
factorised screen relies on---altitude--rigidity separability, the
altitude-independence of solar modulation, and the altitude-dependence of
the logarithmic dose gradient.

%-----------------------------------------------------------------------
\section{Materials and methods}

\subsection{CARI-7 reference dataset}
\label{sec:cari7}
Reference effective-dose rates were generated with CARI-7 (version 4.1.3),
the physics-based GCR model maintained by the FAA Civil Aerospace Medical
Institute \citep{anic2016,copeland2008,faacari7}, whose atmospheric
response functions were computed with MCNPX~2.7.0. All values reported here
are direct CARI-7 executions. We used ICRP Publication~103 effective dose
from all particles (the program's effective-dose tally, consistent with
fluence-to-dose conventions, \citealp{ferrari1997}), the 2004 ISO/Nymmik
GCR local interstellar spectrum, and monthly-average mode; solar modulation
is set by the heliocentric potential (HCP), a heliospheric-modulation proxy
\citep{caballero2004,usoskin2017}, from the program's date-indexed FAA
table \citep{faahcp}. For each case the program returns the vertical cutoff
rigidity ($R_c$) of the location and the effective dose rate.

The references correspond to Korean Peninsula airspace (Incheon,
$37.47^{\circ}$N, $126.45^{\circ}$E), for which CARI-7 returns
$R_{c,\rm ref}=10.52$\,GV (2025 epoch). The nominal solar input was
November~2025, HCP\,$=1{,}141$\,MV (a solar-maximum-phase month of Solar
Cycle~25; the smoothed maximum was October~2024, \citealp{sidc2025}). GCR
dose is \emph{lowest} near solar maximum, so this reference is the lower
bound; a solar-minimum reference (January~2020, HCP\,$=275$\,MV) gives the
conservative upper bound (Section~\ref{sec:annual}). The framework is a
quiet-time GCR baseline and excludes SEP events and ground-level
enhancements, such as GLE~77 (11~November~2025) \citep{gle77}, which
require real-time systems such as NAIRAS and WASAVIES \citep{mertens2013}.

A vertical altitude profile was computed at Incheon from 10{,}000 to
40{,}000\,ft in 5{,}000\,ft steps (seven levels). For the latitude
correction, effective-dose rates at FL300 were computed at eleven
geographic locations spanning polar to equatorial routes
(Section~\ref{sec:flat}), with three additional locations reserved as a
held-out test set. Vertical altitude profiles were additionally computed
at a polar (Svalbard, $R_c=0.03$\,GV) and an equatorial (Singapore,
$R_c=17.2$\,GV) location to test the altitude--rigidity separability
assumption directly (Section~\ref{sec:latext}).

\subsection{Altitude screening formula}
\label{sec:empformula}
Over the operational FL200--FL400 band the altitude--dose relation was fit
by least squares as a second-order polynomial in $\ln D$ versus altitude
$h$ (km), using the five cruise-band CARI-7 nodes. A quadratic (rather than
a single exponential, i.e.\ a linear $\ln D$--$h$ law) is required because
the logarithmic dose gradient decreases steadily with altitude as the
secondary cascade approaches its transition region. Fit quality is reported
as the mean absolute relative error (MARE) against CARI-7.

\subsection{Geomagnetic latitude correction}
\label{sec:flat}
A latitude correction $f_{\rm lat}(R_c)=D(R_c)/D(R_{c,\rm ref})$ was
derived from FL300 effective-dose rates computed directly with CARI-7 at
eleven geographic nodes whose returned vertical cutoff rigidities
\citep{smart2009cutoff} span $R_c=0.0$--17.2\,GV (Table~\ref{tab:nodes}),
normalised to the Incheon
reference ($R_{c,\rm ref}=10.52$\,GV). A saturation model
$f_{\rm lat}(R_c)=(1+aR_{c,\rm ref}^{\,n})/(1+aR_c^{\,n})$ was fitted and
compared with a power law and an exponential decay; parameter uncertainty
was obtained by bootstrap resampling ($B=10{,}000$). The fitted model was
then tested on three held-out nodes (Munich, Rome, Tokyo; $R_c=4.0$, 6.3,
11.3\,GV) that were not used in the fit, giving an out-of-sample error.
Longitudinal variation at fixed $R_c$ (east--west effect, South Atlantic
Anomaly) is a second-order effect and is not modelled; the altitude--$R_c$
separation is tested directly in Section~\ref{sec:latext}.

\subsection{Annual dose and regulatory context}
\label{sec:annual}
Annual effective doses were estimated for four altitude levels by
multiplying CARI-7 mean dose rates by indicative annual flight hours.
Three solar scenarios were used, all from real CARI-7 runs: nominal solar
maximum (HCP\,$=1{,}141$\,MV), solar minimum (HCP\,$=275$\,MV), and a
conservative upper bound (solar minimum $\times1.15$). Projections were
compared with the ICRP public reference level (1\,mSv\,yr$^{-1}$;
\citealp{icrp2007}) and the Korean aircrew management criterion
(6\,mSv\,yr$^{-1}$; \citealp{knssc2021}); reference levels for highly
exposed aircrew are typically 5--10\,mSv\,yr$^{-1}$ \citep{icrp132}. The
1\,mSv\,yr$^{-1}$ value is used here as a conservative screening action
threshold, not as the applicable aircrew dose limit.

\subsection{Consistency with CARI-7 and published literature}
\label{sec:consistency}
We do not claim an independent validation against in-flight measurements.
CARI-7 is itself validated against tissue-equivalent proportional counter
and silicon-detector measurements
\citep{badhwar2002,dachev2010,bennett2013,waker2019,dyer2007} and in
cross-code and in-flight intercomparisons
\citep{eurados2012,bottollier2004,deangelis2009,matthia2014pandoca,meier2018},
which agree to within about $\pm30\%$. The screen is instead
evaluated by how closely it reproduces CARI-7, including an out-of-sample
test on held-out latitude nodes (Section~\ref{sec:flat}). Published
in-flight dose rates at comparable geomagnetic conditions
\citep{lee2015cosmic,spurny2001,matthia2014pandoca} are discussed
qualitatively as context; a strict quantitative comparison is limited by
the conversion between the measured ambient dose equivalent $H^*(10)$ and
the CARI-7 effective dose $E$, whose operational-quantity ratio at aviation
altitudes is close to unity but carries appreciable spread, and by the
reference dosimetry standard in use \citep{icru84,iso207851}. We therefore
treat the measurement literature as supporting the reference rather than as
a direct test of the screening formula.

\subsection{Use of artificial-intelligence tools}
Large-language-model assistance (Anthropic Claude) was used for LaTeX
editing, figure-generation code, and language polishing. All study design,
CARI-7 executions, numerical results, and conclusions are the authors' own
and were checked by the authors; no text or data were accepted without
verification.

%-----------------------------------------------------------------------
\section{Results}

\subsection{Altitude screening formula}
The cruise-band fit (Figure~\ref{fig:alt}) is
\begin{equation}
  \ln D\,(\text{\textmu Sv\,h}^{-1})
  = -0.0232\,h_{\rm km}^{2} + 0.701\,h_{\rm km} - 4.172,
  \label{eq:logquad}
\end{equation}
with $h_{\rm km}$ the altitude in kilometres. Over the FL200--FL400 band
Equation~\ref{eq:logquad} reproduces the CARI-7 reference with
MARE\,$=0.10\%$ (maximum residual $0.2\%$; Table~\ref{tab:ef}); refitting
over the extended FL100--FL400 range gives
$\ln D=-0.0217\,h_{\rm km}^{2}+0.673\,h_{\rm km}-4.059$ with
MARE\,$=1.05\%$. The local logarithmic gradient
$d\ln D/dh$ falls from $\approx0.38$\,km$^{-1}$ near FL200 to
$\approx0.17$\,km$^{-1}$ near FL400, so a single exponential
(constant-gradient) law cannot reproduce the band; the quadratic term
captures this curvature while remaining a closed-form pocket calculation.

\begin{figure}[htbp]
\centering
\includegraphics[width=0.92\textwidth]{figure1.eps}
\caption{Altitude--dose-rate relation from CARI-7 (Incheon,
$R_c=10.52$\,GV, HCP\,$=1{,}141$\,MV). The log-quadratic screen
(Equation~\ref{eq:logquad}) is drawn solid over its calibrated FL200--FL400
band and dotted outside it. Open circles are CARI-7 executions; the
secondary axis is altitude in thousands of feet.}
\label{fig:alt}
\end{figure}

\begin{table}[htbp]
\centering
\caption{CARI-7 reference dose rates and the log-quadratic screen
(Equation~\ref{eq:logquad}) over the FL200--FL400 cruise band (Incheon,
$R_c=10.52$\,GV, HCP\,$=1{,}141$\,MV). RE\,$=$\,(screen$-$CARI-7)/CARI-7.}
\label{tab:ef}
\small
\begin{tabular}{lccccc}
\hline\hline
\textbf{Profile} & \textbf{Alt. (ft)} & \textbf{Alt. (km)} &
\textbf{CARI-7} & \textbf{Screen} & \textbf{RE} \\
 & & & (\textmu Sv\,h$^{-1}$) & (\textmu Sv\,h$^{-1}$) & (\%) \\
\hline
Very high & 40{,}000 & 12.19 & 2.506 & 2.504 & $-0.1$ \\
High      & 35{,}000 & 10.67 & 1.931 & 1.933 & $+0.1$ \\
Reference & 30{,}000 & 9.14  & 1.340 & 1.340 & $+0.0$ \\
Medium    & 25{,}000 & 7.62  & 0.835 & 0.834 & $-0.2$ \\
Low edge  & 20{,}000 & 6.10  & 0.465 & 0.466 & $+0.1$ \\
\hline
\multicolumn{6}{@{}p{0.95\linewidth}@{}}{\footnotesize MARE $=0.10\%$ over
FL200--FL400 (five nodes). Below FL200 the FL100--FL400 refit should be
used (MARE $1.05\%$; see text).}\\
\hline\hline
\end{tabular}
\end{table}

\subsection{Solar-cycle modulation}
CARI-7 runs at the solar-minimum reference (HCP\,$=275$\,MV) give dose
rates a uniform $+15\%$ above the solar-maximum reference across all
altitudes (e.g.\ $+15.2\%$ at both FL300 and FL400), consistent with the
12--20\% full-cycle span reported elsewhere
\citep{mertens2013,clairand2023,dobynde2021} and with heliospheric
modulation theory \citep{caballero2004,usoskin2017}. Because the scaling is
nearly
altitude-independent, it is well represented by a single logarithmic HCP
term, $1-\alpha\ln(\text{HCP}/\text{HCP}_{\rm ref})$ with
$\alpha=0.108$ (Section~\ref{sec:latext}). Altitude is therefore the
dominant predictor at any phase of the cycle, with solar modulation a
second-order correction.

\subsection{Annual doses and regulatory screening}
Operations sustained at $\sim$40{,}000\,ft (672\,h\,yr$^{-1}$) give
1.68--2.23\,mSv\,yr$^{-1}$, exceeding the 1\,mSv\,yr$^{-1}$ public
reference level under all three solar scenarios (Figure~\ref{fig:annual},
Table~\ref{tab:annual}). Standard widebody cruise at $\sim$30{,}000\,ft
(576\,h\,yr$^{-1}$) gives 0.77--1.02\,mSv\,yr$^{-1}$, crossing the
threshold only under the conservative scenario; operations at or below
20{,}000\,ft remain well below all thresholds. These values are consistent
with reported long-haul aircrew doses and national per-capita estimates
\citep{grajewski2011,hammer2009,clairand2023,copeland2008,yasuda2018,hubert2018,zarrouk2008,bennett2013}
and should be read as a screening trigger for closer monitoring, not as
evidence of health effects.

\begin{figure}[htbp]
\centering
\includegraphics[width=0.85\textwidth]{figure3.eps}
\caption{Estimated annual effective dose by altitude and solar scenario,
from CARI-7 dose rates and indicative annual flight hours. Dashed line:
ICRP public reference level (1\,mSv\,yr$^{-1}$).}
\label{fig:annual}
\end{figure}

\begin{table}[htbp]
\centering
\caption{Annual effective dose (mSv\,yr$^{-1}$) by altitude and solar
scenario, from real CARI-7 dose rates. $\dagger$ exceeds the
1\,mSv\,yr$^{-1}$ public reference level; conservative $=$ solar
minimum $\times1.15$.}
\label{tab:annual}
\small
\begin{tabular}{lcccc}
\hline\hline
\textbf{Altitude} & \textbf{h\,yr$^{-1}$} & \textbf{Nominal} &
\textbf{Solar min} & \textbf{Conservative} \\
\hline
$\sim$40{,}000\,ft & 672 & 1.684$^\dagger$ & 1.940$^\dagger$ & 2.231$^\dagger$ \\
$\sim$30{,}000\,ft & 576 & 0.772 & 0.889 & 1.022$^\dagger$ \\
$\sim$20{,}000\,ft & 360 & 0.168 & 0.194 & 0.223 \\
$\sim$10{,}000\,ft & 600 & 0.067 & 0.077 & 0.089 \\
\hline\hline
\end{tabular}
\end{table}

\subsection{Geomagnetic latitude extension}
\label{sec:latext}
Within the Korean mid-latitude envelope altitude dominates, but cutoff
rigidity becomes a first-order correction toward the poles (lower $R_c$,
higher dose) and the equator (higher $R_c$, lower dose). Fitting the
eleven-node correction factor (Table~\ref{tab:nodes}), the saturation
model was preferred ($R^2=0.996$) over a power law and an exponential
decay. Bootstrap resampling ($B=10{,}000$) gave $\hat a=0.025$
(95\% CI 0.012--0.046) and $\hat n=1.50$ (95\% CI 1.27--1.78). The
complete screen combines the altitude base, the latitude factor, and a
logarithmic HCP term:
\begin{equation}
  D(h,R_c,\text{HCP}) = D_{\rm alt}(h)
  \,\frac{1+a R_{c,\rm ref}^{\,\hat n}}{1+a R_c^{\,\hat n}}
  \,\bigl[1-\alpha\ln(\text{HCP}/\text{HCP}_{\rm ref})\bigr],
  \label{eq:full}
\end{equation}
with $D_{\rm alt}(h)$ from Equation~\ref{eq:logquad}, $a=0.024$,
$\hat n=1.52$, and $\alpha=0.108$. On the eleven fitted nodes the latitude
factor has an in-sample MARE of 3.4\%; on the three held-out nodes the
out-of-sample MARE is 1.9\% (Figure~\ref{fig:fig6}), a genuine test of the
fitted form. The correction spans a factor of 2.5 between polar
($f_{\rm lat}=1.81$) and equatorial ($f_{\rm lat}=0.73$) nodes.

Applied at FL300 for 600\,h\,yr$^{-1}$, the fitted model places the
1\,mSv\,yr$^{-1}$ crossover near $R_c\approx7.3$\,GV: estimated annual
doses range from 1.48\,mSv\,yr$^{-1}$ on polar routes ($R_c\to0$) to
0.53\,mSv\,yr$^{-1}$ on equatorial routes ($R_c\approx17$\,GV). The Korean
reference itself (0.80\,mSv\,yr$^{-1}$ at FL300 on this basis) falls below
the threshold, consistent with Section~\ref{sec:annual}.

The screen multiplies a single altitude shape (Equation~\ref{eq:logquad})
by the latitude factor, which assumes the altitude profile is invariant
with cutoff rigidity. CARI-7 altitude profiles computed directly at a
polar ($R_c=0.03$\,GV) and an equatorial ($R_c=17.2$\,GV) location test
this: the equatorial profile matches the reference shape to within
MARE\,$=0.4\%$ over FL200--FL400, confirming separability at mid-to-low
latitudes. Toward the geomagnetic pole the softer primary spectrum makes
the profile steepen, and the separable screen deviates by up to $+9\%$
(FL200) to $-9\%$ (FL400) there (MARE\,$=5.6\%$). The full screen is thus
most accurate for $R_c\gtrsim4$\,GV; polar-route screening carries this
additional $\sim$5--9\% separability uncertainty, which is acceptable for
a first-pass flag but is stated explicitly.

\begin{table}[htbp]
\centering
\caption{FL300 effective-dose rates computed with CARI-7 at eleven
geographic nodes (HCP\,$=1{,}141$\,MV), with the vertical cutoff rigidity
$R_c$ returned by CARI-7 and the resulting
$f_{\rm lat}=D/D_{\rm ref}$ ($D_{\rm ref}=1.340$\,\textmu Sv\,h$^{-1}$ at
$R_{c,\rm ref}=10.52$\,GV).}
\label{tab:nodes}
\scriptsize
\setlength{\tabcolsep}{4pt}
\begin{tabular}{cccl}
\hline\hline
\textbf{$R_c$ (GV)} & \textbf{$D$ (\textmu Sv\,h$^{-1}$)} &
\textbf{$f_{\rm lat}$} & \textbf{Representative location} \\
\hline
0.00 & 2.425 & 1.810 & Resolute Bay (74.7$^\circ$N) \\
0.03 & 2.425 & 1.810 & Svalbard (78.3$^\circ$N) \\
0.48 & 2.425 & 1.810 & Keflav\'ik (64.0$^\circ$N) \\
1.85 & 2.398 & 1.790 & Glasgow (55.9$^\circ$N) \\
3.74 & 2.175 & 1.624 & Paris (49.0$^\circ$N) \\
4.86 & 2.010 & 1.501 & Zagreb (45.7$^\circ$N) \\
8.34 & 1.573 & 1.175 & Athens (37.9$^\circ$N) \\
10.52 & 1.340 & 1.000 & Incheon / Korea (reference) \\
15.91 & 1.024 & 0.764 & Hong Kong (22.3$^\circ$N) \\
16.79 & 0.988 & 0.738 & Manila (14.5$^\circ$N) \\
17.21 & 0.971 & 0.725 & Singapore (1.4$^\circ$N) \\
\hline\hline
\end{tabular}
\end{table}

\begin{figure}[htbp]
\centering
\includegraphics[width=\textwidth]{figure6.eps}
\caption{Geomagnetic latitude extension. (a)~Correction factor
$f_{\rm lat}(R_c)$ from the eleven CARI-7 fit nodes (circles) with the
saturation model; three held-out nodes (triangles) are overplotted.
(b)~Bootstrap distribution of the exponent $\hat n$ ($B=10{,}000$):
$\hat n=1.50$, 95\% CI [1.27, 1.78]. (c)~Screen vs.\ CARI-7 for the
in-sample (MARE 3.4\%) and held-out (MARE 1.9\%) nodes.}
\label{fig:fig6}
\end{figure}

\subsection{Application to operational routes}
\label{sec:routes}
To test the screen as a decision tool, CARI-7 full great-circle doses were
computed for five real routes from Incheon (RKSI) spanning the latitude
range (Table~\ref{tab:routes}). The screen applies Equation~\ref{eq:full}
at the cruise flight level and the great-circle-midpoint cutoff rigidity,
multiplied by block time---inputs a user has without running any code. For the four long-haul routes the hand formula
reproduces the CARI-7 route dose to within $-2\%$ to $+8\%$. The short
domestic sector (ICN--CJU) is overestimated by 54\%, as expected where
climb and descent dominate a cruise-based screen, but its absolute dose
($<1$\,\textmu Sv) is negligible. Most importantly, for a representative
crew schedule the screen and CARI-7 return the \emph{same}
1\,mSv\,yr$^{-1}$ exceedance flag on all five routes: the trans-polar
(ICN--JFK), trans-Siberian (ICN--FRA) and North-Pacific (ICN--LAX)
long-haul rotations are flagged, the equatorial (ICN--SIN) and domestic
(ICN--CJU) sectors are not.

\begin{table}[htbp]
\centering
\caption{CARI-7 full great-circle effective dose for five routes from
Incheon (Nov 2025, HCP\,$=1{,}141$\,MV) and the screen
(Equation~\ref{eq:full}) evaluated at the cruise flight level and the
great-circle-midpoint cutoff rigidity, multiplied by block time.}
\label{tab:routes}
\small
\begin{tabular}{lccccc}
\hline\hline
\textbf{Route} & \textbf{Cruise} & \textbf{$R_{c,\rm mid}$} &
\textbf{Block} & \textbf{Screen} & \textbf{CARI-7} \\
 & (FL) & (GV) & (h) & (\textmu Sv) & (\textmu Sv) \\
\hline
ICN--JFK (polar)       & 350 & 0.01  & 14.2 & 50.4 & 47.3 \\
ICN--FRA (trans-Siberia)& 370 & 1.27  & 11.7 & 45.1 & 42.6 \\
ICN--LAX (N.\ Pacific) & 360 & 2.89  & 11.2 & 37.7 & 35.0 \\
ICN--SIN (equatorial)  & 370 & 16.37 & 6.7  & 10.1 & 10.3 \\
ICN--CJU (domestic)    & 280 & 11.31 & 1.0  & 1.1  & 0.69 \\
\hline
\multicolumn{6}{@{}p{0.95\linewidth}@{}}{\footnotesize Long-haul RE
(screen$-$CARI-7): $+6.6\%$ (JFK), $+5.8\%$ (FRA), $+7.7\%$ (LAX),
$-2.2\%$ (SIN). ICN--CJU $+54\%$ (climb/descent-dominated; dose
negligible). The 1\,mSv\,yr$^{-1}$ annual flag agrees with CARI-7 on all
five routes.}\\
\hline\hline
\end{tabular}
\end{table}

%-----------------------------------------------------------------------
\section{Discussion}

The curvature of $\ln D$ with altitude (Equation~\ref{eq:logquad}) is
physically expected: the logarithmic gradient flattens as the secondary
cascade approaches its transition region \citep{schraube1998}, so
single-exponential relations fitted over particular bands
\citep{vukovic2022,bartlett2004} are band-dependent. The second-order term
reproduces CARI-7 across FL200--FL400 to 0.1\% (FL100--FL400 to
$\sim$1\%) while remaining hand-evaluable.

The contribution is deliberately modest and operational. Full route codes
(CARI-7, EPCARD, AVIDOS, SIEVERT) and global-dose-field representations
\citep{tobiska2018,meier2018,phoenix2026} remain the tools of record for
regulatory dosimetry; what they do not provide is a single closed-form
expression, evaluable without software or network access, for a first-pass
screen of which profiles approach the 1\,mSv\,yr$^{-1}$ action level.
Equations~\ref{eq:logquad} and~\ref{eq:full} supply that screen with a
stated $\le1\%$ reproduction of CARI-7 in altitude and a held-out-tested
latitude factor---to our knowledge the first hand-evaluable screen anchored
to a named reference code and validated out-of-sample across both altitude
and latitude. It is of practical value for crew briefing, retrospective
estimation from flight logs, and resource-limited operations.

Beyond the formula, the CARI-7 runs characterise the structure any
factorised screen depends on, with transferable results. First,
altitude--rigidity separability is bounded: the normalised altitude profile
is invariant to $\le1\%$ for $R_c\gtrsim4$\,GV and steepens by
$\sim$5--9\% only toward the pole, so one altitude shape suffices outside
polar routes. Second, full-cycle solar modulation is altitude-independent
at $\sim$15\%, captured by a one-parameter logarithmic HCP term. Third, the
logarithmic dose gradient falls monotonically with altitude (0.38 to
0.17\,km$^{-1}$), so single-exponential screens are intrinsically
band-limited. These are properties of the aviation dose field, not fit
artefacts.

Several limitations bound the intended use. The multiplicative
altitude--$R_c$ separation carries the $\sim$5--9\% polar uncertainty
quantified in Section~\ref{sec:latext}, and the latitude factor is tested
out-of-sample over only three nodes, so a broader held-out set and
cross-epoch testing would further constrain it. Finally, the framework is a
quiet-time GCR baseline: SEP
events and ground-level enhancements, and storm-time cutoff suppression,
are outside its scope and require real-time systems such as NAIRAS and
WASAVIES \citep{mertens2013,mewaldt2005,kubo2015wasavies,sato2018wasavies,meier2014swindex}.
Extending and testing the screen
against measured in-flight datasets spanning a range of latitudes,
altitudes, and solar conditions would broaden its operational basis.

%-----------------------------------------------------------------------
\section{Conclusion}

A compact log-quadratic in altitude,
$\ln D=-0.0232\,h_{\rm km}^{2}+0.701\,h_{\rm km}-4.172$, reproduces the
CARI-7 cruise-band effective-dose reference to within 0.1\% (1\% over
FL100--FL400) and requires only flight altitude; a saturation-type
cutoff-rigidity factor---fitted to direct CARI-7 runs at eleven nodes and
validated on three held-out nodes to 1.9\%---and a logarithmic HCP term
extend it across latitude and the solar cycle. Used as a screen, the
formula flags sustained high-altitude operations
(1.7--2.2\,mSv\,yr$^{-1}$ near 40{,}000\,ft) and lower-rigidity routes
($R_c\lesssim7.3$\,GV at FL300) as exceeding the 1\,mSv\,yr$^{-1}$ public
reference level. It is intended as a minimum-data complement to full
route-based dosimetry, not as a replacement, and its stated error and
validity range make its screening role explicit.

%-----------------------------------------------------------------------
\section*{Acknowledgments}
This research was conducted as part of a project under the Ministry of
National Defense of the Republic of Korea. Heliocentric-potential data
were obtained from the FAA database derived from Apatity Cosmic Ray
Station neutron measurements \citep{faahcp}; CARI-7 was provided by the
FAA Civil Aerospace Medical Institute.

\section*{Funding}
This research received no external funding.

\section*{Conflicts of interest}
The authors declare no conflicts of interest.

\section*{Data availability statement}
The CARI-7 output files, derived data, and analysis code are archived on
Zenodo (\url{https://doi.org/10.5281/zenodo.19560458}; CC-BY-4.0). All
reference values are reproducible with the FAA CARI-7 tool (v4.1.3) using
the parameters in Section~\ref{sec:cari7}: ICRP~103 effective dose, 2004
ISO/Nymmik GCR spectrum, monthly-average mode, and the stated dates and
coordinates.

\section*{Author contribution statement}
K.\ Kang: conceptualisation, methodology, software, writing.
D.-H.\ Lee and D.\ Lee: supervision, funding acquisition, review.
J.-M.\ Park, H.S.\ Chi, J.\ Yoo, M.\ Shin, H.\ Kim, J.\ Yoo:
investigation, data curation, formal analysis, visualisation. All authors
reviewed and approved the manuscript.

\section*{Ethics approval}
Not applicable (no human or animal subjects).

\section*{Informed consent}
Not applicable.

%-----------------------------------------------------------------------
\bibliographystyle{apalike}
\bibliography{myref}

\end{document}
